\chapter{تابع مهلة تسوية الضبط التمويني حتى قرار الإحالة}
\label{ch:settlement}

ضبط تمويني بحق فرن له مهلة تسوية سبعة أيام. الإجراء يبلّغ ببدء المهلة، وينتظرها، ثم يقرأ الملف من جديد، ويطلب قرار رئيسة الدائرة إن لم تُسوَّ المخالفة، وينفّذ قرارها.

\section{افهم القصة قبل البناء}
\label{sec:real-story}
\begin{rcquick}{افهم القصة قبل البناء}
\qphase{التحضير}
\qwho{أنت}{قبل الدخول إلى كيان}
\qwhere{هذه الخطوة}
\qdo{افتح «البيانات» واقرأ كيف تقابل كل مرحلة حقيقية خطوة في الإجراء.}
\qdo{فرّق بين ثلاثة أشياء: الضبط، والمهلة، وقرار الإحالة.}
\qdo{اسأل المديرية: هل المهلة أيام تقويمية أم أيام عمل؟ وأي المخالفات تجوز تسويتها؟}
\qresult{تعرف ما يفعله الإجراء وما يبقى قراراً لإنسان.}
\qwhy{لبعض المخالفات يجيز قانون حماية المستهلك (المرسوم التشريعي رقم 8 لعام 2021) التسوية خلال سبعة أيام، وإلا يُحال الضبط إلى القضاء. الإجراء يتابع المهلة بدل الذاكرة.}
\qdata{اعتماد الضبط}{يصبح الفرن مخالفاً بانتظار التسوية}{مشغّل «تحديث سجل» مع حارس على «الحالة الرقابية».}
\qdata{بدء المهلة}{يُبلَّغ فادي ومازن}{«إرسال إشعار».}
\qdata{انتظار المهلة}{سبعة أيام}{«تأخير».}
\qdata{الفحص بعد المهلة}{هل سُوّيت؟}{«جلب سجل» ثم «تفرّع شرطي».}
\qdata{قرار الإحالة}{رئيسة الدائرة تحيل أو تؤجل}{«طلب موافقة» بإجراءين مخصصين.}
\qdata{الإحالة}{ملف إحالة ووضع الفرن}{«تحديث سجل كيان» و«إنشاء مهمة».}
\qdata{التدريبي}{الأشخاص والأفران والضبوط}{تحمل «تجريبي» أو «تدريب».}
\end{rcquick}

\section{تأكد من الفريق والحقل الرقابي}
\label{sec:real-team}
\begin{rcquick}{تأكد من الفريق والحقل الرقابي}
\qphase{التحضير}
\qwho{حساب المسؤول}{مدير التطبيق}
\qwhere{الموارد البشرية ← التعيينات}
\qdo{تأكد أن للثلاثة تعيينات «نشط»، وإلا أنشئ المناصب والأعضاء كما في المثال الأول.}
\qdo{«ادوات البناء» ← «الكيانات» ← «المنشآت التجارية»: تأكد من حقل «الحالة الرقابية» بخياراته الأربعة.}
\qdo{دوّن مفتاح الحقل و«القيمة المخزّنة» لكل خيار.}
\qresult{الحقل ظاهر في ملف «فرن الأمانة» بخياراته الأربعة، وقيمه المخزّنة مدوّنة عندك.}
\qwhy{الإجراء يكتب في حقل ويرسل إلى أشخاص. إن نقص خيار أو تعيين، يتوقف في منتصفه. والشروط تقارن القيمة المخزّنة، لا النص الظاهر.}
\qdata{مازن قاسم — تجريبي}{رئيس دورية حماية المستهلك — 01}{\en{DEMO-CP-004}؛ يُبلَّغ ببدء المهلة.}
\qdata{سلمى عبود — تجريبي}{رئيس دائرة حماية المستهلك}{\en{DEMO-CP-006}؛ تقرر الإحالة أو تأجيلها.}
\qdata{فادي الخوري — تجريبي}{موظف التسويات والإحالة — 01}{\en{DEMO-CP-007}؛ يتابع ويسجّل التسوية.}
\qdata{مفتاح الحقل}{\en{inspection\_status}}{استبدله بمفتاح مؤسستك.}
\qdata{سليمة}{القيمة المخزّنة عندك}{}
\qdata{مخالفة — بانتظار التسوية}{\en{pending\_settlement}}{}
\qdata{سُوّيت المخالفة}{\en{settled}}{}
\qdata{أُحيل الضبط إلى القضاء}{\en{referred}}{}
\end{rcquick}

\section{صمّم القاعدة قبل البناء}
\label{sec:real-design}
\begin{rcquick}{صمّم القاعدة قبل البناء}
\qphase{التحضير}
\qwho{أنت}{قبل البناء}
\qwhere{ورقة جانبية}
\qdo{اقرأ القاعدة في «البيانات» واكتبها بكلماتك.}
\qdo{ارسمها خطوات: مشغّل، حارس، إشعار، تأخير، جلب، تفرّع، طلب، تنفيذ.}
\qdo{أجب: لماذا نقرأ الملف بعد التأخير بدل بيانات المشغّل؟}
\qresult{القاعدة مكتوبة، وتعرف الجواب: ربما سجّل فادي التسوية خلال المهلة.}
\qwhy{الخطأ في قاعدة فيها انتظار أسبوع لا يظهر إلا بعد أسبوع. التصميم الدقيق قبل البناء يوفّر هذا الأسبوع.}
\qdoc{القاعدة}{عندما تتغير «الحالة الرقابية» لفرن إلى «مخالفة — بانتظار التسوية»، أبلغ فادي ومازن. انتظر سبعة أيام. ثم اقرأ ملف الفرن من جديد. إذا لم تعد الحالة «بانتظار التسوية»، فانتهِ. وإلا اطلب من سلمى قرار الإحالة. إذا قررت الإحالة، فاكتب «أُحيل الضبط إلى القضاء» وأنشئ لفادي مهمة ملف الإحالة. وإذا أجّلت، فأبلغ فادي بسببها.}
\end{rcquick}

\section{أنشئ الإجراء بمشغّل «تحديث سجل»}
\label{sec:real-trigger}
\begin{rcquick}{أنشئ الإجراء بمشغّل «تحديث سجل»}
\qphase{البناء}
\qwho{حساب المسؤول}{من يبني الإجراء}
\qwhere{الإجراءات الداخلية ← حماية المستهلك — تدريب ← إجراء جديد}
\qdo{«الاسم»: من «البيانات».}
\qdo{«مزيد من المُشغِّلات» ← ابحث عن «تحديث سجل» ← «إنشاء وفتح».}
\qdo{«الجدول»: «المنشآت التجارية».}
\qresult{المشغّل على اللوحة يعرض «المنشآت التجارية».}
\qwhy{الحالة تتغير بتعديل ملف موجود، لا بإنشائه. لكن «تحديث سجل» يعمل مع كل تعديل، حتى تصحيح هاتف، لذلك نضع حارساً بعده.}
\qdata{الاسم}{متابعة مهلة تسوية الضبط — تدريب}{\req{}}
\qdata{\en{\$event.data}}{الملف بعد التعديل}{}
\qdata{\en{\$event.oldData}}{الملف قبل التعديل}{}
\qdata{\en{\$event.changedFields}}{الحقول التي تغيّرت}{}
\end{rcquick}

\section{ضع الحارس}
\label{sec:real-guard}
\begin{rcquick}{ضع الحارس}
\qphase{البناء}
\qwho{حساب المسؤول}{من يبني الإجراء}
\qwhere{بناء ← «+» بعد المشغّل}
\qdo{أضف «تفرّع شرطي»، والاسم المرجعي \en{guard}.}
\qdo{المسار الأول «بدأت المهلة»: «حدّد الشرط» ← «شرط مخصص». أدرج الحقل بـ«إدراج بيانات» من «المُشغِّل»، وأكمل كما في «البيانات».}
\qdo{على المسار البديل أضف «إنهاء» بالنتيجة \en{not\_relevant}.}
\qresult{باختبار الخطوة: من «سليمة» إلى «بانتظار التسوية» يسلك «بدأت المهلة»؛ وبقاؤها «بانتظار التسوية» يسلك البديل.}
\qwhy{بلا حارس، كل تعديل على فرن مخالف يبدأ مهلة جديدة. الحارس يسمح فقط عندما تنتقل الحالة الآن، بمقارنة \en{data} و\en{oldData}.}
\qdata{الاسم المرجعي}{\en{guard}}{\rec{}}
\qdata{المسار البديل}{إنهاء}{النتيجة \en{not\_relevant}.}
\qdoc{شرط «بدأت المهلة»}{\en{\$event.data.inspection\_status == "pending\_settlement" and \$event.oldData.inspection\_status != "pending\_settlement"} \\ \en{and} تعني «و»: يجب أن يصح الجزءان.}
\end{rcquick}

\section{أبلغ فادي ومازن ببدء المهلة}
\label{sec:real-notify}
\begin{rcquick}{أبلغ فادي ومازن ببدء المهلة}
\qphase{البناء}
\qwho{حساب المسؤول}{من يبني الإجراء}
\qwhere{بناء ← «+» على مسار «بدأت المهلة»}
\qdo{أضف «إرسال إشعار». «المستلمون»: فادي ومازن.}
\qdo{العنوان والمحتوى من «البيانات»، واسم الفرن بـ«إدراج بيانات» من «المُشغِّل».}
\qdo{«متقدم»: الاسم المرجعي \en{notify\_start}، وفعّل «إعادة المحاولة قبل الاستسلام».}
\qresult{الخطوة «جاهز»، والعنوان يحمل مرجع اسم الفرن.}
\qwhy{يجب أن يعرف موظف التسويات أن العدّ بدأ اليوم، ويعرف رئيس الدورية أن ضبطه دخل المتابعة.}
\qdata{عنوان الرسالة}{بدأت مهلة تسوية: + اسم الفرن}{}
\qdata{محتوى الرسالة}{اعتُمد ضبط بحق هذه المنشأة. مهلة التسوية سبعة أيام من اليوم. سجّل التسوية في الملف إن تمت.}{}
\qdata{الاسم المرجعي}{\en{notify\_start}}{}
\end{rcquick}

\section{انتظر المهلة بخطوة «تأخير»}
\label{sec:real-delay}
\begin{rcquick}{انتظر المهلة بخطوة «تأخير»}
\qphase{البناء}
\qwho{حساب المسؤول}{من يبني الإجراء}
\qwhere{بناء ← «+» ← تدفق التحكم ← تأخير}
\qdo{أضف «تأخير» بعد الإشعار.}
\qdo{اختر الوضع والمدة بحسب سياسة المديرية.}
\qdo{الاسم المرجعي \en{wait\_period}.}
\qresult{تعرف أن العملية ستبقى سبعة أيام «بانتظار استجابة» عند هذه الخطوة، وأن هذا ليس خطأً.}
\qwhy{«تأخير» يوقف هذا التشغيل وحده ثم يكمل من حيث توقف. لكل ضبط عملية تشغيل بمهلتها، ولا يتأثر فرن بآخر.}
\qdata{الوضع}{«لمدة محددة» أو «أيام العمل»}{\req{} «أيام العمل» تتخطى العطل بحسب تقويم المؤسسة.}
\qdata{المدة}{\en{7 days}، أو 7 أيام عمل}{\req{}}
\qdata{في الساعة}{09:00}{\opt{} مع «أيام العمل».}
\qdata{الاسم المرجعي}{\en{wait\_period}}{\rec{}}
\end{rcquick}

\section{اقرأ ملف الفرن من جديد}
\label{sec:real-get}
\begin{rcquick}{اقرأ ملف الفرن من جديد}
\qphase{البناء}
\qwho{حساب المسؤول}{من يبني الإجراء}
\qwhere{بناء ← «+» ← السجلات ← جلب سجل}
\qdo{«الجدول» «المنشآت التجارية»، و«السجل» بتعبير \en{\$event.data.id}.}
\qdo{«إذا كان معرّف السجل فارغاً»: «أفشِل الخطوة».}
\qdo{الاسم المرجعي \en{recheck}. «متقدم» ← «اتباع مسار الفشل».}
\qdo{على مسار الفشل: إشعار لسلمى بالنص من «البيانات».}
\qresult{يخرج من \en{recheck} خطان؛ حذف الملف أثناء المهلة لن يُسقط المتابعة بصمت.}
\qwhy{\en{\$event.data} صورة الملف لحظة الضبط. إن سجّل فادي التسوية في اليوم الثالث، فالصورة القديمة لا تعرف. «جلب سجل» يقرأ الملف الآن.}
\qdata{الاسم المرجعي}{\en{recheck}}{\req{} يستخدمه التفرّع التالي.}
\qdata{مسار الفشل}{إشعار لسلمى}{تعذّرت قراءة ملف منشأة بعد مهلة التسوية؛ ربما حُذف. راجع الضبط يدوياً.}
\end{rcquick}

\section{تفرّع: هل ما زال بانتظار التسوية؟}
\label{sec:real-branch}
\begin{rcquick}{تفرّع: هل ما زال بانتظار التسوية؟}
\qphase{البناء}
\qwho{حساب المسؤول}{من يبني الإجراء}
\qwhere{بناء ← «+» بعد \en{recheck}}
\qdo{أضف «تفرّع شرطي»، الاسم المرجعي \en{still\_pending}.}
\qdo{المسار «لم تُسوَّ»: «شرط مخصص» ← «إدراج بيانات» ← خطوة \en{recheck} ← «الحالة الرقابية»، وأكمل بـ\en{== "pending\_settlement"}.}
\qdo{المسار البديل: «إنهاء» بالنتيجة \en{settled\_or\_changed}.}
\qresult{باختبار التفرّع: «بانتظار التسوية» تسلك «لم تُسوَّ»، و«سُوّيت المخالفة» تسلك البديل.}
\qwhy{الطلب يذهب إلى سلمى فقط إن لم تتم التسوية. وإلا ينتهي التشغيل بنتيجة واضحة يقرؤها من يراجع العملية.}
\qdata{حقول \en{recheck} لا تظهر}{شغّل «اختبار هذه الخطوة» على \en{recheck} أولاً}{القائمة تعرض حقول خطوة عملت في اختبار.}
\qdoc{شرط «لم تُسوَّ»}{\en{\$steps.recheck.output.<مفتاح الحالة> == "pending\_settlement"} \\ أدرج الجزء الأول بـ«إدراج بيانات» ولا تكتبه بيدك.}
\end{rcquick}

\section{اطلب قرار سلمى بأزرار مخصصة}
\label{sec:real-approval}
\begin{rcquick}{اطلب قرار سلمى بأزرار مخصصة}
\qphase{البناء}
\qwho{حساب المسؤول}{من يبني الإجراء}
\qwhere{بناء ← «+» على مسار «لم تُسوَّ» ← طلب موافقة}
\qdo{املأ العنوان والتفاصيل والمرحلة والمهلة من «البيانات».}
\qdo{«إجراءات الطلب»: الإجراء الأول للإحالة، والثاني للتأجيل، بالقيم في «البيانات».}
\qdo{«مرفقات المكتبة»: زر التعبير ← \en{[]}.}
\qdo{الاسم المرجعي \en{referral}.}
\qresult{لا مشكلات على الخطوة في المراجعة، وتحت «مرفقات المكتبة» \en{[]}.}
\qwhy{سلمى تقرر «إحالة إلى القضاء» أو «تأجيل الإحالة»، لا «موافقة». نتيجة كل زر هي ما يقرؤه التفرّع بعده.}
\qdata{عنوان الموافقة}{انتهت مهلة التسوية: + اسم الفرن من \en{recheck}}{\req{}}
\qdata{تفاصيل الموافقة}{انقضت سبعة أيام على اعتماد الضبط ولم تُسجَّل تسوية. قرّري الإحالة إلى القضاء أو تأجيلها مع ذكر السبب.}{\rec{}}
\qdata{مراحل الموافقة}{مرحلة واحدة؛ سلمى}{\req{}}
\qdata{المهلة الزمنية}{\en{3 days}}{\rec{}}
\qdata{إرفاق أحدث إصدار من السجل}{مفعّل؛ \en{\$event.data.id}}{\rec{} ترى سلمى الملف الذي تقرر عليه.}
\qdata{الإجراء 1}{\en{refer}؛ «إحالة إلى القضاء»؛ \en{referred}}{«إكمال»؛ التوقيع «مطلوب»؛ التعليق «اختياري».}
\qdata{الإجراء 2}{\en{postpone}؛ «تأجيل الإحالة»؛ \en{postponed}}{«إكمال»؛ التوقيع «لا شيء»؛ التعليق «مطلوب».}
\qdata{مرفقات المكتبة}{\en{[]}}{\req{} وإلا «Internal command failed with status 422».}
\qdata{تسمية أو نتيجة فارغة}{«مفتاح النتيجة مطلوب.»}{لا تترك أيّاً منهما.}
\end{rcquick}

\section{نفّذ القرار: إحالة أو تأجيل}
\label{sec:real-outcome}
\begin{rcquick}{نفّذ القرار: إحالة أو تأجيل}
\qphase{البناء}
\qwho{حساب المسؤول}{من يبني الإجراء}
\qwhere{بناء ← «+» بعد \en{referral} ← تفرّع شرطي}
\qdo{ثلاثة مسارات: «إحالة» و«تأجيل» بالشرطين في «البيانات»، والبديل.}
\qdo{«إحالة»: «تحديث سجل كيان» (\en{\$event.data.id}، «أُحيل الضبط إلى القضاء») ثم «إنشاء مهمة» لفادي.}
\qdo{«تأجيل»: إشعار لفادي ثم «إنهاء» بـ\en{postponed}.}
\qdo{البديل (انتهت مهلة سلمى دون رد): إشعار لسلمى وفادي ثم «إنهاء».}
\qresult{«كل الخطوات جاهزة»، ولكل مسار نهاية.}
\qwhy{قرار سلمى يجب أن يظهر في ملف الفرن وعمل فادي. ولا تجعل الإجراء يحيل وحده: القرار لصاحب الصلاحية، والإجراء يذكّر ويجهّز.}
\qdata{عنوان المهمة}{إعداد ملف إحالة: + اسم الفرن}{موعد بعد يومي عمل.}
\qdata{إشعار التأجيل}{أجّلت رئيسة الدائرة الإحالة. افتح الطلب لقراءة السبب.}{}
\qdata{إشعار البديل}{انتهت مهلة قرار الإحالة دون رد.}{}
\qdoc{شرطا التفرّع}{\en{\$steps.referral.output.terminalOutcomeKey == "referred"} \\ \en{\$steps.referral.output.terminalOutcomeKey == "postponed"}}
\end{rcquick}

\section{اختبر المسارات قبل النشر}
\label{sec:real-test}
\begin{rcquick}{اختبر المسارات قبل النشر}
\qphase{الاختبار والتشغيل}
\qwho{حساب المسؤول}{من يبني الإجراء}
\qwhere{المحرر ← اختبار}
\qdo{في «بيانات الاختبار» اكتب كائناً فيه \en{data} و\en{oldData} لكل من الحالتين الأوليين، ثم «محاكاة (آمنة)».}
\qdo{للحالات الثلاث الأخيرة: «اختبار هذه الخطوة» على \en{still\_pending} والتفرّع بعد \en{referral}.}
\qdo{قارن كل نتيجة بالجدول في «البيانات».}
\qresult{طابقت كل حالة نتيجتها. إن سلك تعديل الهاتف طريق المهلة، فالحارس لا يقارن \en{oldData}.}
\qwhy{هذا الإجراء ينتظر أسبوعاً قبل أن يكشف خطأه. المحاكاة واختبار الخطوات يكشفان الشرط الخاطئ اليوم.}
\qdata{تعديل هاتف؛ الحالة قبل وبعد «بانتظار التسوية»}{ينتهي عند الحارس بـ\en{not\_relevant}}{لا إشعار.}
\qdata{من «سليمة» إلى «بانتظار التسوية»}{إشعار، ثم يقف عند «تأخير»}{}
\qdata{\en{recheck} يعيد «سُوّيت المخالفة»}{ينتهي بـ\en{settled\_or\_changed}}{لا طلب لسلمى.}
\qdata{\en{recheck} «بانتظار التسوية» ونتيجة \en{referred}}{تحديث الحالة ومهمة فادي}{}
\qdata{نتيجة \en{postponed}}{إشعار فادي ثم \en{postponed}}{}
\qdoc{بيانات اختبار: بدأت المهلة}{\en{\{"data": \{"id": "test-1", "inspection\_status": "pending\_settlement"\}, "oldData": \{"id": "test-1", "inspection\_status": "<قيمة سليمة>"\}\}}}
\end{rcquick}

\section{انشر وطبّق على «فرن الأمانة»}
\label{sec:real-live}
\begin{rcquick}{انشر وطبّق على «فرن الأمانة»}
\qphase{الاختبار والتشغيل}
\qwho{سلمى، ثم فادي}{رئيس دائرة حماية المستهلك؛ موظف التسويات والإحالة — 01}
\qwhere{مراجعة ونشر ← المنشآت التجارية ← فرن الأمانة}
\qdo{المسؤول: «مراجعة ونشر» ← «نشر \en{v1}».}
\qdo{سلمى: اعتمد ضبطاً تدريبياً، أو غيّر «الحالة الرقابية» من «سليمة» إلى «مخالفة — بانتظار التسوية».}
\qdo{فادي: الإشعارات ← «بدأت مهلة تسوية: فرن الأمانة».}
\qdo{المسؤول: «عمليات التشغيل» ← «فعلي». ثم عدّل هاتف الفرن.}
\qresult{عملية واحدة «بانتظار استجابة» عند «تأخير» مع موعد الاستئناف. تعديل الهاتف لم يبدأ عملية ثانية.}
\qwhy{هنا ترى الإجراء مع عمل حقيقي: يبدأ العدّ وحده لحظة تغيّر الحالة، ويتجاهل التعديلات الأخرى.}
\qdata{للتدريب: لا تنتظر أسبوعاً}{انسخ الإجراء في مسودة تدريبية بمدة \en{5 minutes}}{انشرها وحدها ثم أوقفها؛ لا تغيّر مدة الإجراء الحقيقي.}
\end{rcquick}

\section{سجّل تسوية خلال المهلة}
\label{sec:real-settle}
\begin{rcquick}{سجّل تسوية خلال المهلة}
\qphase{الاختبار والتشغيل}
\qwho{فادي}{موظف التسويات والإحالة — 01}
\qwhere{المنشآت التجارية ← فرن الأمانة}
\qdo{غيّر «الحالة الرقابية» إلى «سُوّيت المخالفة» واحفظ.}
\qdo{بعد انتهاء المدة، افتح العملية في «عمليات التشغيل».}
\qdo{اضغط خطوة \en{recheck} واقرأ ما أعادته.}
\qresult{\en{recheck} قرأت «سُوّيت المخالفة»، وانتهت العملية بـ\en{settled\_or\_changed}، ولم يصل سلمى طلب.}
\qwhy{أهم دليل على صحة القاعدة أن تبقى صامتة حين يكون العمل قد تم.}
\end{rcquick}

\section{جرّب طريق الإحالة حتى نهايته}
\label{sec:real-refer}
\begin{rcquick}{جرّب طريق الإحالة حتى نهايته}
\qphase{الاختبار والتشغيل}
\qwho{سلمى، ثم فادي}{رئيس دائرة حماية المستهلك؛ موظف التسويات والإحالة — 01}
\qwhere{طلبات بانتظارك}
\qdo{أنشئ ضبطاً تدريبياً على «فرن الجامعة — تدريب» ودعه بلا تسوية حتى تنتهي المدة.}
\qdo{سلمى: «طلبات بانتظارك» ← «انتهت مهلة التسوية: فرن الجامعة — تدريب» ← «إحالة إلى القضاء» ← وقّع.}
\qdo{فادي: افتح مهامه وملف الفرن.}
\qdo{كرر بفرن آخر مع «تأجيل الإحالة».}
\qresult{«أُحيل الضبط إلى القضاء»، ومهمة «إعداد ملف إحالة: فرن الجامعة — تدريب»، والعملية «مكتمل». مع التأجيل: التعليق إلزامي، ويصل فادي إشعار، ولا تتغير الحالة.}
\qwhy{طريق الإحالة له أثر على الملف وعمل فادي، فيجب أن يُرى مرة حقيقية قبل الاعتماد عليه.}
\end{rcquick}

\section{تابع الإجراء أسبوعياً}
\label{sec:real-monitor}
\begin{rcquick}{تابع الإجراء أسبوعياً}
\qphase{الاختبار والتشغيل}
\qwho{حساب المسؤول}{من يبني الإجراء}
\qwhere{الإجراءات الداخلية}
\qdo{التصفية «تحتاج إلى متابعة». إن ظهر الإجراء: «فتح التشغيل الفاشل».}
\qdo{الإجراء ← «عمليات التشغيل» ← «فعلي»: العمليات عند «تأخير» ضبوط في المهلة، وعند \en{referral} قرارات تنتظر سلمى.}
\qdo{إن تغيّرت سياسة المهلة: عدّل «تأخير» وانشر إصداراً جديداً.}
\qresult{تعرف من «عمليات التشغيل» وحدها: كم فرناً في المهلة، وكم قراراً معلقاً، ومتى تنتهي أقرب مهلة.}
\qwhy{الضبوط التي بدأت مهلتها تكمل على إصدارها القديم بعد النشر، وهذا صحيح لأن مهلتها حُسبت عند بدئها.}
\end{rcquick}

\section{طبّق وحدك}
\label{sec:real-practice}
\begin{rcquick}{طبّق وحدك}
\qphase{الإتقان}
\qwho{أنت}{دون الرجوع إلى الخطوات}
\qwhere{الإجراءات الداخلية}
\qdo{أضف «إرسال تذكير بالموافقة» لسلمى إذا لم ترد خلال يوم، واختبره.}
\qdo{ابنِ إجراءً بمشغّل «جدولة» يرسل لسلمى كل أحد عدد الأفران «بانتظار التسوية» بـ«البحث عن سجلات كيان».}
\qdo{اشرح لزميل الفرق بين بيانات المشغّل ونتيجة «جلب سجل».}
\qresult{نفّذت تمريناً واحداً على الأقل وحدك، واختبرته بالمحاكاة ثم فعلياً.}
\qwhy{الإجراء المجدول يكمل ما تعلمته: لا ينتظر حدثاً على سجل، بل يبدأ في موعد ويبحث عن السجلات بنفسه.}
\end{rcquick}

\section{اختبر فهمك}
\label{sec:real-hero}
\begin{rcquick}{اختبر فهمك}
\qphase{الإتقان}
\qwho{أنت}{أجب قبل أن تفتح «البيانات»}
\qwhere{الإجراء وعملياته}
\qdo{أجب عن كل سؤال بصوت عالٍ.}
\qdo{افتح «البيانات» وقارن.}
\qdo{تحقق من كل جواب في المكان المذكور.}
\qresult{تشرح لماذا سلك كل تشغيل طريقه، من التطبيق وحده.}
\qwhy{تصبح مرجع فريقك في الإجراءات عندما تشرح السبب، لا الأزرار فقط.}
\qdata{لماذا لم يبدأ تعديل الهاتف مهلة جديدة؟}{الحارس قارن \en{oldData} فسلك البديل}{النتيجة \en{not\_relevant} في العملية.}
\qdata{متى تنتهي مهلة «فرن الأمانة»؟}{موعد الاستئناف}{خطوة «تأخير» في العملية المنتظرة.}
\qdata{لماذا لم يصل سلمى طلب عن فرن سُوّي؟}{\en{recheck} قرأت الحالة الجديدة}{مخرجات \en{recheck} والمسار البديل.}
\qdata{من قرر إحالة «فرن الجامعة» ومتى؟}{سلمى، في ردها على \en{referral}}{ردود الخطوة، والطلب في «مكتملة» عندها.}
\qdata{لماذا لم يتأثر ضبط قديم بتغيير المدة؟}{العمليات الجارية تبقى على إصدارها}{رقم الإصدار في العملية و«سجل الإصدارات».}
\qdata{ما الفرق بين \en{\$event.data} و\en{\$steps.recheck.output}؟}{الأول لحظة المشغّل، والثاني لحظة القراءة}{«إدراج بيانات».}
\qdata{لماذا لا يحيل الإجراء وحده؟}{القرار لصاحب الصلاحية}{خطوة «طلب موافقة».}
\end{rcquick}
